\documentclass[12pt]{article}
\usepackage[utf8]{inputenc}
\usepackage[T1]{fontenc}
\usepackage[margin=1in]{geometry}
\usepackage{setspace}
\usepackage{xcolor}
\usepackage[colorlinks=true, urlcolor=blue!60!black, linkcolor=blue!60!black]{hyperref}
\setlength{\parindent}{2em}
\setlength{\parskip}{1.25em}
\renewcommand{\baselinestretch}{1.0}
\begin{document}
\begin{center}
{\LARGE \textit{Shurtleff v. City of Boston}, 596 U.S. 243 (2022)\par}
\vspace{0.8em}
{\large Case Brief --- Four-Page Edition\par}
\end{center}
\vspace{1em}

\section{Memory Jogger}

Boston had allowed private groups to raise about 50 different flags at 284 City Hall ceremonies without any recorded refusal but refused Camp Constitution's Christian flag; the Court held 9--0 that, because Boston did not actively control the flags' messages, the program was private speech, and the refusal was unconstitutional viewpoint discrimination (\href{https://supreme.justia.com/cases/federal/us/596/20-1800/}{596 U.S. 243}).

\section{Detailed Case Facts}

Three 83-foot flagpoles stand on City Hall Plaza, a 7-acre space Boston acknowledged was a public forum. The first flies the United States flag, the second the Massachusetts flag, and the third ordinarily the city flag. Since at least 2005, Boston let private groups hold ceremonies on the plaza and temporarily raise a flag of their choosing on the third pole. Between 2005 and 2017 it approved 284 consecutive ceremonies involving about 50 flags---mostly foreign national flags, but also the Pride Flag, an emergency medical services banner, and a local credit union's flag. It had no written content policy and had never rejected an application.

In July 2017, Harold Shurtleff, director of Camp Constitution, asked to hold a ceremony to ``commemorate the civic and social contributions of the Christian community'' and raise what he called the ``Christian flag.'' Commissioner Gregory T. Rooney denied the request, the first denial in the program's history, solely because he feared flying a religious flag would violate the Establishment Clause, and offered to allow the ceremony with a different flag.

\section{Procedural History}

Shurtleff and Camp Constitution sued Boston and Rooney under the Free Speech Clause. After denying preliminary relief, the District Court granted summary judgment for Boston on a stipulated record, holding the displays government speech, and the First Circuit affirmed. The Supreme Court granted certiorari, heard argument on January 18, 2022, and reversed on May 2, 2022.

\section{Judicial Votes}

Reversed and remanded, \textbf{9--0}.
\begin{itemize}
\item \textbf{Majority:} Justice Breyer, joined by Chief Justice Roberts and Justices Sotomayor, Kagan, Kavanaugh, and Barrett.
\item \textbf{Concurrence:} Justice Kavanaugh.
\item \textbf{Concurring in the judgment:} Justice Alito, joined by Justices Thomas and Gorsuch; Justice Gorsuch, joined by Justice Thomas.
\item \textbf{Dissents:} None.
\end{itemize}

\section{Holding}

Applying a holistic, context-driven inquiry into history, public perception, and the extent of active government control, the Court held that Boston's flag-raising program was not government speech. Because Boston opened the flagpole to private expression, refusing Camp Constitution's flag because of its religious viewpoint violated the Free Speech Clause.

\section{Analysis of Opinions}

\textbf{Majority (Breyer).} When government speaks for itself, it need not be viewpoint neutral; the ballot box checks its messages (\textit{Summum}; \textit{Walker}). When it invites private participation, courts weigh three factors. \textit{History} favored Boston: flags at seats of government have long conveyed official messages. \textit{Public perception} was inconclusive, since the city lowered its own flag for private ceremonies that passersby might attribute to the groups. \textit{Active control} was decisive: asked how far Boston shaped the flags' messages, the answer was ``not at all.'' Boston told the public it sought ``to accommodate all applicants'' at its ``public forums,'' never inspected flags, had no guidelines, and approved 284 straight requests---a ``come-one-come-all attitude'' like trademark registration in \textit{Matal v. Tam}, unlike \textit{Summum}'s selective monuments or \textit{Walker}'s reviewed plates. Boston could change its policy, as San Jose had with a closed list. With government speech gone, the Establishment Clause defense fell too, and refusing a flag because it promoted a specific religion was viewpoint discrimination under \textit{Good News Club} and \textit{Rosenberger}.

\textbf{Kavanaugh (concurring).} The dispute arose from an official's mistaken belief that treating religious speech equally violates the Establishment Clause; excluding religious speakers from neutral forums violates the First Amendment.

\textbf{Alito (concurring in the judgment).} The three-factor test can mask censorship: heavy control is the hallmark of a royal censor, and perception lets government exploit misattribution. Speech is the government's only if it purposefully communicates a governmentally determined message through authorized agents without abridging private speech; Boston's 284 unvetted flags failed that test.

\textbf{Gorsuch (concurring in the judgment).} Boston's error traced to \textit{Lemon} and its ``reasonable observer'' test, a ``children's game'' that pushes localities to censor religious speech; courts should rely on historical practices and understandings, under which neutral forum access is no establishment.

\section{Examples: Future Applications}

\begin{enumerate}
\item \textbf{Same-Side (Exclusion Invalid): Library Heritage Wall.} Despite a written disclaimer, a library that approves 140 consecutive plaques without review hosts private speech and may not reject a mosque's plaque for its religious viewpoint.
\item \textbf{Same-Side (Exclusion Invalid): Town-Green Holiday Displays.} A town that approves every seasonal display may not refuse an atheist group's solstice sign as offensive.
\item \textbf{Opposite-Side (Government Speech): Closed-List Flag Policy.} A city that adopts a written policy flying only flags chosen by council resolution and closes the pole to private groups is speaking for itself.
\item \textbf{Opposite-Side (Government Speech): Reviewed License Plates.} A board that reviews and owns every specialty plate design controls the message, as in \textit{Walker}, and may reject a slogan.
\item \textbf{Fence-Sitter: Juried Mural Program.} A commission that reviews murals for quality and approves about 90 percent, with no content guidelines, sits between \textit{Summum} and \textit{Shurtleff}.
\end{enumerate}

\section{Critique}

\textbf{The ``leveling down'' paradox.} By confirming that Boston could escape viewpoint neutrality with a written closed-list policy, \textit{Shurtleff} encourages cities to shut public forums altogether, reducing private speech to avoid controversial flags such as the Satanic Temple's later application in Boston.

\textbf{Control as sword and shield.} Alito's concurrence exposes a tension in Breyer's third factor: the more aggressively a city polices and rejects private speech, the more likely a court is to reward that censorship by calling the program government speech.

\section{Key Quotations}

\begin{quote}
``In answering these questions, we conduct a holistic inquiry designed to determine whether the government intends to speak for itself or to regulate private expression. Our review is not mechanical; it is driven by a case's context rather than the rote application of rigid factors.''
\end{quote}
\noindent --- Justice Breyer, Opinion of the Court, 596 U.S. at 252.

\begin{quote}
``All told, while the historical practice of flag flying at government buildings favors Boston, the city's lack of meaningful involvement in the selection of flags or the crafting of their messages leads us to classify the flag raisings as private, not government, speech---though nothing prevents Boston from changing its policies going forward.''
\end{quote}
\noindent --- Justice Breyer, Opinion of the Court, 596 U.S. at 258.

\begin{quote}
``To avoid a spurious First Amendment problem, Boston wound up inviting a real one. Call it a \textit{Lemon} trade\ldots This Court long ago interred \textit{Lemon}, and it is past time for local officials and lower courts to let it lie.''
\end{quote}
\noindent --- Justice Gorsuch, concurring in the judgment, 596 U.S. at 280, 288.

\end{document}
